% !PRINTED 307
\weberpage{307}

\gesperrt{Es wird also weder ein $VP$ noch ein $PP$ verloren oder gewonnen.}

Wenn aber $f(x)$ selbst durch Null geht, so haben wir, wenn wir $f_1(x)$ positiv voraussetzen, was mit Rücksicht auf die vorher gemachte Bemerkung genügt,
\[
\begin{array}{c|cc}
 & 0 & 1 \\ \hline
f(x-0) & - & + \\
f(x+0) & + & + \\ \hline
F & + & +
\end{array}
\text{.}
\]
% [resolved 2026-08-06, 5x zoom strip 2/5] "in" is correct. At zoom the apparent ascender is
% a speck of ink above the n, not an h, and sense settles it independently: the clause reads
% "Es geht also hier ein VP in ein PP ueber", where "ih" is not a word. Transcription stands.
Es geht also hier ein $VP$ in ein $PP$ über, d. h. \gesperrt{es wird ein $VP$ verloren, ein $PP$ gewonnen.}

Wenn endlich $F_\nu$ durch Null geht, wo natürlich $0 < \nu < n$ sein muss, so ist das Verhalten folgendes:

Wenn $F_\nu = 0$ ist, so müssen nothwendig $f_{\nu-1}$ und $f_{\nu+1}$ gleiche Vorzeichen haben, da sonst $F_\nu$ aus zwei positiven Gliedern bestehen würde, und nicht verschwinden könnte. Es ist ferner $F_\nu'(x)$ nach (6) vom selben Vorzeichen wie das Product $f_{\nu-1} f_\nu F_{\nu+1}$, und also hat $F_\nu(x-0)$ das entgegengesetzte, $F_\nu(x+0)$ das gleiche Vorzeichen, wie das Product.

Auch hier können wir in der Aufzählung die Fälle weglassen, die aus den aufgeführten durch gleichzeitige Zeichenänderung in allen drei Functionen $F_{\nu-1}, F_\nu, F_{\nu+1}$ entstehen, und es bleiben also folgende Fälle:
% [sic] The first right-hand header is printed with a prime after $\nu+1$; the corresponding later header is unprimed.
\[
\begin{array}{c|ccc||ccc}
f & \nu-1, & \nu, & \nu+1 & \nu-1, & \nu, & \nu+1' \\
 & + & + & + & + & + & + \\ \hline
F(x-0) & + & - & + & - & - & + \\ \hline
F(x+0) & + & + & + & - & + & +
\end{array}
\]
\[
\begin{array}{c|ccc||ccc}
f & \nu-1, & \nu, & \nu+1 & \nu-1, & \nu, & \nu+1 \\
 & + & - & + & + & - & + \\ \hline
F(x-0) & + & + & + & - & + & + \\ \hline
F(x+0) & + & - & + & - & - & +
\end{array}
\text{.}
\]
Die Zeichenwechsel sind in diesen vier Fällen:
%CONT>>
